\naslov{Predstavljivost števil}

\newcommand{\nan}{\text{nan}}

Množica predstavljivih števil $P(b, t, L, U)$ vsebuje vsa števila oblike
\[
  \pm 0, c_1 c_2 \ldots c_t \cdot b^e,
\]
kjer je $L \le e \le U$ in $c_i \in \{0, \ldots, b-1\}$.
Dovoljena so tudi števila $0$, $\infty$, $-\infty$ ter $\nan$.
Števila so lahko normalizirana ali denormalizirana; pri prvem tipu velja $c_1
\ne 0$.
Denormalizirana števila so dovoljena le, če se tega števila ne mora zapisati v
normalizirani obliki.
Običajen \verb+float+ je razreda $P(2, 24, -125, 128)$, \verb+double+ pa $P(2,
53, -1021, 1023)$.

Najbližje predstavljivo število številu $x$ v neki množici zapišemo s $\fl(x)$.
Napako operacij lahko izračunamo kot \pojem{absolutno napako} $d_a = \hat{x} -
x$ in \pojem{relativno napako} $d_r = d_a / x$.
Pri ocenjevanju lahko uporabimo dejstvo
\[
  \frac{\abs{\fl{x} - x}}{\abs{x}} \le u = \pol b^{1-t}.
\]
Standard IEEE zagotavlja omejeno napako pri osnovnih operacijah,
\begin{itemize}
\item $\fl{x \oplus y} = (x \oplus y)(1+\delta)$ za $\abs{\delta} \le u$ in
  osnovne operacije $+, -, \cdot, /$,
\item $\fl{\sqrt{x}} = \sqrt{x} (1+\delta)$ za $\abs{\delta} \le u$.
\end{itemize}

\naslov{Nelinearne enačbe}

Če iščemo ničlo funkcije $f(x)$ na nekem intervalu, lahko uporabimo bisekcijo
ali katero od iteracijskih metod.
V splošnem je problem za večkratne ničle bolj občutljiv, za $m$-kratno ničlo v
splošnem velja
\[
  \abs{x - \alpha} \le {\left( \frac{m!\,\varepsilon}{\abs{f^{(m)}(\alpha)}}
  \right)}^{1/m},
\]
če je $\abs{f(x)} \le \varepsilon$.

Pri navadni iteraciji iščemo fiksno točko dane funkcije $g$ pri nekem začetnem
približku, s postopkom $x_{r+1} = g(x_r)$.
Da bo postopek konvergiral, mora biti $g$ v okolici fiksne točke skrčitev.
Dovolj je že, da je $\abs{g'(\alpha)} < 1$ in $g$ zvezno odvedljiva.
Če vemo, da je $g$ Lipschitzova s konstanto $L < 1$, lahko ocenimo
\[
  \abs{x_r - \alpha} \le L^r \abs{x_0 - \alpha}
\]
oziroma
\[
  \abs{x_{r+1} - \alpha} \le \frac{L}{1 - L} \abs{x_{r+1} - x_r},
\]
kjer je $\alpha$ fiksna točka.
Interval konvergence lahko torej določamo s pomočjo odvoda (izračunamo, kje je
$\abs{g'} < 1$), vendar nam to ne da vedno celotnega intervala.
Drug način je, da direktno dokažemo, kakšna je limita zaporedja $(x_r)_r$.
Vizualno si lahko pomagamo s pajčevinastim diagramom.

Pri konvergenci poznamo tudi \pojem{red}.
To je število $p$, za katerega velja
\[
  \lim_{r \to \infty} \frac{\abs{x_{r+1} - \alpha}}{\abs{x_r - \alpha}^{p}} = C
\]
za neko konstanto $C > 0$.
Zadosten pogoj za red $p$ je, da velja $g'(\alpha) = \cdots = g^{(p-1)}(\alpha)
= 0$ ter $g^{(p)} \ne 0$.
Višji kot je red, boljša je metoda.

Če iščemo ničlo funkcije $f$ in želimo zagotoviti red $2$ (za enostavno ničlo),
lahko uporabimo \pojem{tangentno metodo}
\[
  g(x) = x - \frac{f(x)}{f'(x)}.
\]
Če je $f$ vsaj dvakrat odvedljiva in $\alpha$ ničla kratnosti $m$, potem za $g$
iz tangentne metode velja
\[
  \lim_{x \to \alpha} g'(x) = 1 - \inv{m}.
\]
Podobna je \pojem{sekantna metoda}, kjer imamo predpis
\[
  x_{r+1} = x_r - \frac{f(x_r) (x_r - x_{r-1})}{f(x_r) - f(x_{r-1})}.
\]
To pa sedaj ni navadna iteracija.

Če imamo dano neko iteracijo
\[
  x_r = F(x_{r-1}, x_{r-2}, \ldots, x_{r-k}),
\]
in vemo, da postopek konvergira (npr.~če dokažemo, da je zaporedje omejeno in
monotono), potem lahko limito izračunamo s pomočjo trika
\[
  \lim_{r \to \infty} x_r = \lim_{r \to \infty} x_{r+i},
\]
torej je limita $\alpha = F(\alpha, \alpha, \ldots, \alpha)$.

\naslov{Matrične norme}

\begin{definicija}
  Preslikava $\norm{\cdot} : \C^{n \times n} \to \R$ je \pojem{matrična norma},
  če velja
  \begin{itemize}
  \item $\norm{A} \ge 0$ in $\norm{A} = 0 \nt A = 0$,
  \item $\norm{\alpha A} = \abs{\alpha} \norm{A}$,
  \item $\norm{A + B} \le \norm{A} + \norm{B}$,
  \item $\norm{A B} \le \norm{A} \norm{B}$.
  \end{itemize}
\end{definicija}

Vsota absolutnih vrednosti vseh elementov matrike $N_1(A)$ in maksimum
absolutnih vrednosti vseh elementov matrike $N_\infty(A)$ žal nista normi,
2-norma vektorizirane matrike pa dejansko je;
\[
  \norm{A}_F = \sqrt{\sum_{i,j} \abs{a_{ij}}^2}.
\]
Velja $\norm{A}_F^2 = \sled (A^H A)$.
Poleg tega je matrična norma tudi $\norm{A} = n N_\infty(A)$, kjer je $n$
dimenzija matrike.
Najbolj uporabne pa so operatorske norme; če je $\norm{\cdot}_v$ vektorska
norma, je pripadajoča matrična norma
\[
  \norm{A}_m = \max_{x \ne 0} \frac{\norm{Ax}_v}{\norm{x}_v}.
\]
Pri tem je operatorska norma $\norm{A}_1$ enaka največji 1-normi stolpca matrike
$A$, $\norm{A}_\infty$ pa največji 1-normi vrstice.
Najbolj uporabna in najtežje izračunljiva je 2-norma, ki je enaka največji
singularni vrednosti matrike (t.j.~korenu lastne vrednosti matrike $A^H A$).
Za poljubno matrično normo in poljubno lastno vrednost $\lambda$ matrike $A$
velja $\abs{\lambda} \le \norm{A}$.

Računanje druge norme je snov predmeta NLA, lahko pa jo ocenimo z enim od
\begin{gather*}
  \inv{\sqrt{n}} \norm{A}_F \le \norm{A}_2 \le \norm{A}_F \\
  \inv{\sqrt{n}} \norm{A}_1 \le \norm{A}_2 \le \sqrt{n} \norm{A}_1 \\
  \inv{\sqrt{n}} \norm{A}_\infty \le \norm{A}_2 \le \sqrt{n} \norm{A}_\infty \\
  \norm{A}_2 \le \sqrt{\norm{A}_1 \norm{A}_\infty} \\
  N_\infty(A) \le \norm{A}_2 \le n N_\infty(A)
\end{gather*}
Za Frobeniusovo in drugo normo dodatno velja $\norm{A}_F = \norm{A^H}_F$ oziroma
$\norm{A}_2 = \norm{A^H}_2$.

\naslov{LU razcep}

Če imamo dano nesingularno matriko $A$, jo lahko razcepimo v produkt $A = LU$,
kjer je $L$ spodnje diagonalna z enicami na diagonali, in $U$ zgornje diagonalna
matrika.
Pri postopku definiramo matrike $L_k = I - l_k e_k^T$ z inverzom $L_k^{-1} = I +
l_k e_k^T$, kjer je vektor $l_k$ definiran kot
\[
  l_k =
  \begin{bmatrix}
	0 \\ \vdots \\ 0 \\ 1 \\ \nicefrac{w_{k+1}}{w_k} \\ \vdots \\ \nicefrac{w_n}{w_k}
  \end{bmatrix}.
\]
Pri tem je $w$ vektor, ki mu $L_k$ uniči vse komponente pod $k$-to.
Matrika $L_{n-1} \ldots L_2 L_1 A = U$ je zgornje trikotne oblike.
Potem je $L$ kar $L = L_1^{-1} L_2^{-1} \ldots L_{n-1}^{-1}$.
Dobimo ga lahko kar tako, da zapišemo vektorje $l_k$ kot stolpce v matriko.
Ker vemo, da ima $L$ na diagonali enice, jih lahko izpustimo, in zapišemo
celoten razcep kar v eno matriko.
Potem lahko sistem $Ax = b$ rešimo v dveh korakih z enostavno premo oz.~obratno
substitucijo $L y = b$, $Ux = y$.
Enoličen $A = LU$ razcep obstaja natanko tedaj, ko so vse vodilne podmatrike $A$
nesingularne.

Bolj zapomnljivo lahko postopek opišemo na sledeč način.
Ko smo na $k$-tem koraku, razdelimo matriko $A$ na šest delov;
\[
  A =
  \begin{bmatrix}
	A_{11} & A_{12} & A_{13} \\
	A_{21} & A_{22} & A_{23}
  \end{bmatrix}.
\]
Pri tem je $A_{11}$ matrike dimenzij $(k-1) \times (k-1)$, $A_{12}$ in $A_{22}$
pa sta stolpca.
Korake štejemo od $1$ naprej.
Elemente podmatrik $A_{11}$, $A_{12}$, $A_{13}$ in $A_{21}$ samo prepišemo.
Elemente podmatrike $A_{22}$ delimo z diagonalnim elementom $a_{kk}$ (najnižjim
v matriki $A_{12}$), za elemente matrike $A_{23}$ pa naredimo naslednje:
Če želimo transformirati element $a_{ij}$, vzamemo elementa $x$ in $y$, ki sta
direktno levo in direktno nad elementom $a_{ij}$ v blokih $\tilde{A}_{22}$
oziroma $A_{13}$ v novi (tisti, ki je sedaj na pol generirana) matriki, in
popravimo $\tilde{a}_{ij} \gets a_{ij} - xy$.
Na koncu dobimo LU razcep v kompresirani obliki.

Ves ta opis je deloval za LU razcep brez pivotiranja.
Če izvajamo delno pivotiranje, permutiramo neobravnavane vrstice tako, da na
mesto pivota pride po absolutni vrednosti čim večji element.
Tako dobimo razcep $PA = LU$, kjer je $P$ permutacijska matrika.
Pri ročnem postopku to naredimo tako, da menjamo vrstice v spodnjem delu
razdeljene matrike, preden izvedemo korak.
Permutacijsko matriko $P$ potem dobimo tako, da začnemo z identiteto, in na njej
izvajamo enake menjave vrstic kot v postopku.
Podobno teče LU razcep s kompletnim pivotiranjem, kjer si dovolimo tudi menjavo
stolpcev.
Dobimo razcep $PAQ = LU$, kjer matriki $P$ in $Q$ izračunamo posebej; pri obeh
začnemo z identiteto, med postopkom pa na $P$ izvajamo le menjave vrstic, na $Q$
pa le menjave stolpcev.

\naslov{Razcep Choleskega}

Matrika $A$ je simetrična pozitivno definitna natanko tedaj, ko zanjo obstaja
razcep Choleskega $A = V V^T$, kjer je $V$ spodnje trikotna matrika s
pozitivnimi diagonalnimi elementi.
Razcep Choleskega poiščemo z enostavnim postopkom; za vsak $k$ med $1$ in $n$
nastavimo
\[
  v_{kk} = \sqrt{a_{kk} - \sum_{i=1}^{k-1} v_{ki}^2},
\]
nato pa še za $j = k+1, \ldots, n$
\[
  v_{jk} = \inv{v_{kk}} \left( a_{jk} - \sum_{i=1}^{k-1} v_{ki} v_{ji} \right).
\]
V praksi to izgleda tako, da od diagonalnega elementa odštejemo kvadrat vsakega
elementa levo od njega, in dobljeno korenimo; nato pa elementom pod njem
odštejemo skalarni produkt $k$-te in $j$-te vrstice (do $k$-tega stolpca),
in razliko delimo z $v_{kk}$.

\naslov{Sistemi nelinearnih enačb}

Reševanja sistema nelinearnih enačb se lahko lotimo z navadno iteracijo.
Če je $G = (G_1, \ldots, G_n)$ iteracijska funkcija, poznamo dva načina
iteracije:
\begin{itemize}
\item Jacobijeva iteracija: $x^{[r+1]} = G(x^{[r]})$
\item Seidlova iteracija: $x^{[r+1]} = (G_1(x^{[r]}), G_2(x_1^{[r+1]},
  x_2^{[r]}, \ldots, x_n^{[r]}), \ldots, G_n(x_1^{[r+1]}, \ldots,
  x_{n-1}^{[r+1]}, x_n^{[r]}))$
\end{itemize}
Jacobijeva iteracija konvergira na množici, kjer je $G$ skrčitev, za Seidlovo
metodo pa velja podobno za preslikavo $H$, ki jo dobimo tako, da v poznejše
argumente $G$ vstavljamo prejšnje slike.
Zadosten pogoj je, da je največja absolutna lastna vrednost matrike $DG$ v
poljubni točki omejena z nekim $m < 1$, ali pa da je $\norm{JG(x)} < 1$ v
poljubni matrični normi.

Podobno kot v drugem poglavju imamo tu Newtonovo metodo
\[
  G(x) = x - (DF(x))^{-1} F(x),
\]
kjer pa v praksi namesto računanja inverza rešimo linearen sistem.

\naslov{Predoločeni sistemi}

Rešitev predoločenega sistema lahko dobimo z reševanjem normalnega sistema
\[
  A^T A x = A^T b,
\]
kjer je matrika $A^T A$ simetrična pozitivno definitna, torej lahko sistem
rešimo z razcepom Choleskega.

Če iščemo minimum $\norm{A x - b}_2$ pod pogojem $Cx = d$, kjer je prvi sistem
predoločen in drugi nedoločen, rešitev dobimo kot $\tilde{x}$ v
\[
  \begin{bmatrix}
	A^T A & C^T \\ C & 0
  \end{bmatrix}
  \cdot
  \begin{bmatrix}
	\tilde{x} \\ \tilde{y}
  \end{bmatrix}
  =
  \begin{bmatrix}
	A^T b \\ d
  \end{bmatrix}.
\]

% LocalWords:  denormalizirana bisekcijo iteracijskih konvergiral kratnosti LU
% LocalWords:  vektorizirane enice podmatrik pivotiranja iteracijska
% LocalWords:  Predoločeni
